\section{The domestic clause}
\label{sec:6.5}
The term has a second clause, drawn not from the Charter but from the belligerent's own constitution. A machine takes no part in a use of force by a state whose constitution requires legislative authorization for it, unless that authorization exists. For the United States that means a declaration of war or an authorization for the use of military force that covers the conflict. The War Powers Resolution's sixty days \autocite{wpr1973}, run without one, is evidence the authorization is absent; a concurrent resolution directing withdrawal, adopted by both chambers, is evidence of the same kind, though since Chadha it binds nothing, and a joint resolution passed and vetoed is stronger, because the President had to act to defeat it. The clause reads markers, and the markers take a reading. That reading is a legal judgment in a separation-of-powers dispute, and a vendor that publishes it has taken a side in one. So the clause is a term for a statute, not for a licence: it binds through the only route that binds every vendor at once, adoption in peacetime by a legislature with its own reason to give its authorization teeth. The Charter clause can be a licence term because the markers its corroborated licence reads, a Council authorization or a court's judgment, take no reading; its emergency licence reads the defender's claim, and the cases below show how much reading that takes. Until the statute exists, what a vendor can hold on this ground is a refusal plus a record. Whether a declaration or an authorization exists is a fact about the statute book. Whether an existing authorization covers this conflict, whether the sixty days have run, and whether a ceasefire stops the clock are legal judgments the executive will dispute, and the clause takes the reading the executive disputes. It doesn't ask the machine to decide what the Constitution permits, which courts decline to decide. It asks whoever writes the deployment's floor to publish the reading, with the markers it rests on, beside the deployment's translation, where it can be contested through the floor court like any other term. The executive's claim that its own action is authorized is an interested claim and is discounted as one. The clause rests on the foundation the rest of the term rests on. A war the legislature never authorized subjects a population abroad to a decision that the state's own people, through their representatives, had no standing to contest either. And it is enactable where the Charter clause is: in peacetime, by a legislature with its own reason to want a procurement term that gives its authorization teeth.

\begin{esbox}
\textbf{The record, as of 24 September 2026.} The 2026 campaign supplies every marker the clause reads, and none was met. No declaration was made and no authorization passed; one was introduced in the House on 7 May and never left committee \autocite{hjres176}. The administration reported to Congress under the War Powers Resolution on 2 March, which started the sixty days. On 1 May, the sixtieth day, the President wrote to Congress that the hostilities begun on 28 February had terminated, resting on a ceasefire declared on 7 April, while a naval blockade continued \autocite{cbs2026terminated,lawfare2026sixtydays}. Hostilities resumed in July under no new authorization. In June both chambers adopted a concurrent resolution directing the President to remove forces from hostilities against Iran, and nothing changed \autocite{hconres86,asil2026iranwpr}. The clause doesn't need the votes. What it reads is the absence of an authorization, and by 1 May that absence was on the record in the President's own letter, which conceded that hostilities had begun and claimed only that they had ended. Appendix~\ref{sec:F.6} gives the record vote by vote.
\end{esbox}
